% Generated draft: requires mathematical and source review.
\documentclass[11pt]{article}
\usepackage[margin=25mm]{geometry}
\usepackage{fontspec}
\setmainfont{Latin Modern Roman}
\usepackage{amsmath,amssymb,mathtools,unicode-math}
\setmathfont{Latin Modern Math}
\usepackage{enumitem,needspace,xurl,hyperref}
\setlength{\parindent}{0pt}
\setlength{\parskip}{5pt}
\setlength{\emergencystretch}{3em}
\newcommand{\N}{\mathbb{N}}
\newcommand{\Z}{\mathbb{Z}}
\newcommand{\Q}{\mathbb{Q}}
\newcommand{\R}{\mathbb{R}}
\newcommand{\C}{\mathbb{C}}
\newcommand{\F}{\mathbb{F}}
\newcommand{\A}{\mathbb{A}}
\newcommand{\PP}{\mathbb{P}}
\newcommand{\E}{\mathbb{E}}
\newcommand{\eps}{\varepsilon}
\newcommand{\dd}{\,\mathrm{d}}
\begin{document}
\section*{UnsolvedMath ID: 800990}
\textbf{Source problem code:} catalogue entry
\subsection{Definitions and mathematical statement}
Let \(\omega=\{0,1,2,\ldots\}\) and let \(2^\omega=\{\alpha:\omega\to\{0,1\}\}\) be Cantor space. The supplied record asks whether the following assertion holds:
\begin{quote}
For every \(X\subseteq 2^\omega\) with \(\mu(X)=1\), there exists \(Y\subseteq 2^\omega\) with \(\mu(Y)=1\) such that, for every chain \(B\subseteq Y\), there is some \(\alpha\in X\) satisfying
\[
[B]=[\alpha].
\]
\end{quote}
The supplied record does not define \(\mu\), the relation with respect to which \(B\) is a chain, or the meanings of \([B]\) and \([\alpha]\). In particular, it does not specify whether the measure conditions require measurable sets or refer to outer measure, nor whether empty chains are included. These unresolved qualifications are retained: the assertion above reproduces the source question but cannot yet be interpreted as a fully specified mathematical statement.
\subsection{Short English statement}
The question attributed to Levin asks whether every set \(X\subseteq 2^\omega\) of measure one admits a set \(Y\subseteq 2^\omega\) of measure one such that each chain in \(Y\) has its bracketed value represented by an element of \(X\). The measure, chain relation, and bracket notation require clarification from the source.
\subsection{Sources}
\begin{itemize}
\item Supplied local UnsolvedMath registry record: permanent ID \(800990\), source problem code \texttt{catalogue entry}, title \emph{Open Problems in Mathematical Logic --- Problem 2 (Levin)}. The registry labels the problem open; this is a supplied status, not an independently verified assessment.
\item Barmpalias, Gao et al., \emph{Open Problems in Mathematical Logic} (2026), Problem 2, \url{https://arxiv.org/abs/2608.26628}. The supplied background attributes the problem/context to Levin; the registry lists Barmpalias, Gao et al. in its proposer field and 2026 in its proposed-year field. These entries do not establish an original proposal date for Levin.
\item The supplied provenance records access on 2026-10-16 and extraction by \texttt{arxiv\_latex}. These are supplied provenance claims; the cited source has not been independently accessed or verified for this draft. Upstream record ID: \(800990\); upstream revision: \texttt{d371ee917ac9647fbb06cd6fc529deec6f3dda82}.
\end{itemize}
\end{document}
